\chapter{Engenharia de Software para Ecossistemas Cognitivos}

\begin{quote}
\textit{``A arquitetura de software é sobre as coisas que são difíceis de mudar depois.''} --- Martin Fowler
\end{quote}

\section{Arquiteturas Orientadas a Eventos}

\subsection{Fundamentos: Event Sourcing, CQRS e Message Brokers}

\textbf{Event-Driven Architecture} (EDA) é o paradigma onde componentes comunicam-se através de eventos imutáveis. O OpenCode implementa EDA via \textbf{MetaBus} (Capítulo 9): um barramento pub/sub onde agentes publicam eventos e inscrevem-se em tópicos.

\textbf{Event Sourcing} (Fowler, 2005): o estado do sistema é reconstruído a partir da sequência de eventos, não armazenado como snapshot. O MetaBus implementa isso via \texttt{EVENTS\_FILE} (\texttt{.mci\_state/metabus\_events.jsonl}) — um log append-only.

\subsection{Padrão Blackboard: Do Hearsay-II ao OpenCode}

O \textbf{padrão Blackboard} (Erman et al., 1980; Nii, 1986) \cite{erman1980hearsay} foi introduzido no sistema Hearsay-II para reconhecimento de fala. A arquitetura consiste em:

\begin{enumerate}
    \item \textbf{Blackboard}: espaço de dados compartilhado (no OpenCode: \texttt{blackboard.tasks});
    \item \textbf{Knowledge Sources} (KS): módulos especializados independentes (128+ agentes);
    \item \textbf{Control Component}: scheduler que decide qual KS ativar (Attention Router).
\end{enumerate}

\begin{figure}[htbp]
\centering
\begin{tikzpicture}[node distance=1.8cm, auto]
    \node[draw, rectangle, minimum width=8cm, minimum height=5cm, fill=gray!5] (bb) at (0,0) {};
    \node[font=\small\bfseries] at (0,2.2) {BLACKBOARD};
    \node[font=\footnotesize] at (0,1.5) {tasks, agent cards, contexto};
    \node[font=\footnotesize] at (0,0.8) {CFP (Call for Proposals)};
    \node[font=\footnotesize] at (0,0.0) {resultados, reflexões};
    
    \node[draw, rectangle, minimum width=2cm, above left=1.5cm of bb] (ks1) {Agente A};
    \node[draw, rectangle, minimum width=2cm, above=1.5cm of bb] (ks2) {Agente B};
    \node[draw, rectangle, minimum width=2cm, above right=1.5cm of bb] (ks3) {Agente C};
    \node[draw, rectangle, minimum width=3cm, below=1.5cm of bb, fill=black!10] (ctrl) {Attention Router};
    
    \draw[<->] (ks1) -- (bb);
    \draw[<->] (ks2) -- (bb);
    \draw[<->] (ks3) -- (bb);
    \draw[->, thick] (ctrl) -- (bb);
\end{tikzpicture}
\caption{Arquitetura Blackboard do OpenCode: agentes (KS) leem/escrevem no quadro negro; Attention Router escalona.}
\label{fig:blackboard}
\end{figure}

\section{Protocolos A2A (Agent-to-Agent)}

O protocolo A2A (Claro et al., 2025) \cite{a2a2025} padroniza a comunicação entre agentes. Cada agente expõe um \textbf{Agent Card}:

\begin{lstlisting}[language=Python, caption={Estrutura de um Agent Card (A2A Protocol)}]
class AgentCard:
    """Identidade e capacidades de um agente no ecossistema."""
    def __init__(self, agent_id, name, description, capabilities, schema):
        self.agent_id = agent_id
        self.name = name
        self.description = description
        self.capabilities = capabilities     # ['coding', 'review', 'math', ...]
        self.input_schema = schema           # JSON Schema esperado
        self.status = "available"            # available, busy, offline
        self.confidence_score = 0.5          # inicial, atualizado pelo Trust Engine
\end{lstlisting}

\subsection{O Ciclo CFP (Call for Proposals)}

O fluxo de delegação no Blackboard segue um leilão descentralizado:

\begin{enumerate}
    \item \textbf{Post}: orquestrador publica task no Blackboard (\texttt{task.post});
    \item \textbf{Match}: Blackboard identifica agentes elegíveis (capacidades compatíveis);
    \item \textbf{CFP}: Blackboard emite Call for Proposals para agentes elegíveis;
    \item \textbf{Volunteer}: agentes voluntariam-se (\texttt{task.volunteer});
    \item \textbf{Assign}: Attention Router seleciona vencedor por Trust Score (maior confiança);
    \item \textbf{Execute}: agente executa (SDD/TDD);
    \item \textbf{Verify}: BehavioralGate avalia resultado;
    \item \textbf{Reflect}: MetaBus registra reflexão, atualiza Confidence Ledger.
\end{enumerate}

\subsection{PRÁXIS 5.1 — Participando do Blackboard como Agente}

\begin{lstlisting}[language=Python, caption={Agente registrando-se e voluntariando-se no Blackboard}]
from mci.metabus import metabus
from mci.blackboard import blackboard, AgentCard

# 1. Registrar agente
card = AgentCard(
    agent_id="coder_python_007",
    name="Python Coder v7",
    description="Especialista em Python com TDD",
    capabilities=["coding", "python", "tdd", "refactoring"],
    schema={"type": "object", "properties": {"code": {"type": "string"}}}
)
metabus.publish("agent.register", card.to_dict(), source_agent="coder_python_007")

# 2. Ouvir CFPs
def on_cfp(event):
    task = event['payload']
    print(f"CFP recebido: {task['description']}")
    # Voluntaria-se se for elegível
    metabus.publish("task.volunteer", {
        "task_id": task['task_id'],
        "agent_id": "coder_python_007"
    }, source_agent="coder_python_007")

metabus.subscribe("task.cfp", on_cfp)
print("Agente registrado e ouvindo CFPs...")
\end{lstlisting}

\section{CI/CD para Agentes Inteligentes}

\subsection{Pipeline de Deploy para Agentes}

Agentes, como qualquer software, beneficiam-se de CI/CD. O OpenCode recomenda:

\begin{itemize}
    \item \textbf{Testes automatizados}: TDD com \texttt{pytest} (Capítulo 8);
    \item \textbf{Deploy canário}: novo agente recebe 5\% do tráfego; se Trust Score mantiver-se, escala para 100\%;
    \item \textbf{Rollback automático}: se BehavioralGate detectar degradação, reverte para versão anterior;
    \item \textbf{Shadow mode}: agente candidato processa em paralelo (sem afetar produção) para avaliação segura.
\end{itemize}

\section{Observabilidade e Debugging}

\subsection{As Quatro Dimensões da Observabilidade}

\begin{enumerate}
    \item \textbf{Métricas}: Trust Score, latência, taxa de sucesso, taxa de slashing — expostas via Prometheus/Grafana;
    \item \textbf{Logs}: MetaBus events (\texttt{.jsonl}) fornecem log estruturado de toda atividade;
    \item \textbf{Tracing}: cada task tem UUID rastreável do post à reflexão (provenance tracking);
    \item \textbf{Alerting}: BehavioralGate dispara alertas quando Trust Score cai abaixo de threshold (0.3).
\end{enumerate}

\begin{thebibliography}{99}
\bibitem{erman1980hearsay} Erman, L. D., et al. (1980). The Hearsay-II Speech-Understanding System. \textit{ACM Computing Surveys}, 12(2), 213--253. \doiurl{10.1145/356810.356816}
\bibitem{nii1986blackboard} Nii, H. P. (1986). Blackboard Systems. \textit{AI Magazine}, 7(3), 38--53.
\bibitem{fowler2005event} Fowler, M. (2005). Event Sourcing. \url{https://martinfowler.com/eaaDev/EventSourcing.html}
\bibitem{kleppmann2017designing} Kleppmann, M. (2017). \textit{Designing Data-Intensive Applications}. O'Reilly Media.
\bibitem{gelernter1985linda} Gelernter, D. (1985). Generative Communication in Linda. \textit{ACM TOPLAS}, 7(1), 80--112. \doiurl{10.1145/2363.2433}
\bibitem{a2a2025} Claro, M., et al. (2025). Agent-to-Agent Protocol with Metacognitive Confidence Scoring. \textit{arXiv:2505.02279}. \doiurl{10.48550/arXiv.2505.02279}
\bibitem{humble2010continuous} Humble, J., \& Farley, D. (2010). \textit{Continuous Delivery}. Addison-Wesley.
\bibitem{beck2003tdd} Beck, K. (2003). \textit{Test-Driven Development: By Example}. Addison-Wesley.
\end{thebibliography}
